\clearpage
\fancyhead[L]{\small ZWEITER COMET VON 1805.}
\hypertarget{native-275-lower}{}
\begin{center}
ZWEITER COMET VON 1805.

\rule{0.6\linewidth}{0.4pt}

Monatliche Correspondenz. Band XIV. S. 75–86. 1806 Juli.

\rule{0.6\linewidth}{0.4pt}

\textsc{Gauss an von Zach.}
\end{center}
\begin{flushright}
Braunschweig, den 20. Mai 1806.
\end{flushright}

\par\indent --- --- --- Es ist schon geraume Zeit, dass Sie mir die Marseiller Original-Beobachtungen der beiden im vorigen Jahre erschienenen Cometen handschriftlich mitzutheilen die Güte hatten. Ich nahm gleich nach Empfang derselben den zweiten von neuem in Rechnung, und hätte Ihnen die Resultate davon also schon vor zwei Monaten schicken können. Inzwischen ist es nicht Nachlässigkeit gewesen, dass ich dies bisher noch nicht gethan habe, sondern die Ursache davon war, dass ich mir vorgenommen hatte, meiner Untersuchung eine grössere Ausdehnung zu geben, und gerade diese Arbeit wurde durch verschiedene andere Abhaltungen verzögert und unterbrochen. Jetzt nehme ich indess nicht länger Anstand, Ihnen meine bisherigen Untersuchungen vorzulegen, die, wenn sie gleich nicht zu dem Ziele geführt haben, das ich anfangs wünschte und hoffte, doch auch so schon zu manchen nicht uninteressanten Betrachtungen Anlass geben werden. Nach diesen Präliminarien will ich nun von meinen Rechnungen der Reihe nach Bericht erstatten.

\par\indent Die Beobachtungen des Hrn.~\textsc{Thulis} sind an sich alle sehr gut, allein seine an der parallactischen Maschine beobachteten Stern-Positionen sind grösstentheils sehr fehlerhaft, besonders die geraden Aufsteigungen. So ist z.~B. die vom 29. November um 25$^{\mathrm{s}}$ in Zeit zu klein, vom 6. December um 30$^{\mathrm{s}}$, und vom 8. December um 38$^{\mathrm{s}}$ zu klein; auch die Declinationen müssen zum Theil um mehrere Minuten geändert werden. Daher ist es kein Wunder, dass vielen Sternen die Anmerkung beigesetzt worden ist, sie seien nicht aufzufinden, ob ich gleich bei weitem den grössten Theil habe reduciren können. Ich habe sämmtliche Beobachtungen neu reducirt, ausgenommen die vom 11. und 22. November (wozu ich
\SourceLeaf{276}%
die Sterne nirgends gefunden habe) und die vom 5. December, wo der Stern zwar bei \textsc{Bode} (64 Piscium)
aber nirgends in der Hist.\ cél.\ vorkommt. Ich habe also zur Nachweisung der Sterne noch folgende
Zusätze zu machen:

\begin{enumerate}
\item[1)] Die Sterne vom 15. und 16. November sind identisch. Hist.\ cél.\ pag.\ 477 (die Conn.\ des tems an
XI besitze ich nicht).
\item[2)] Der Stern vom 23. November steht in der Hist.\ cél.\ pag.\ 20.
\item[3)] Der Stern am 3. December ist 113 Piscium \textsc{Bode}, und steht in der Hist.\ cél.\ pag.\ 37 und 200.
\item[4)] Am 7. December ist 374 Aquarii \textsc{Bode}, 981 \textsc{Mayer}.
\item[5)] Am 8. December ist 359 Aquarii \textsc{Bode}, derselbe, mit welchem Dr.\ \textsc{Olbers} und ich an diesem Abend
den Cometen verglichen haben.
\end{enumerate}

{\rightskip=0pt plus 1em
\par\indent Nachdem ich nunmehr die Beobachtungen sorgfältig reducirt hatte, verbesserte ich darnach meine
parabolischen Elemente. Es fand sich bei dieser Arbeit, dass sich die drei zum Grunde gelegten Örter
in der Parabel nicht ganz so genau, als man wol hätte wünschen können, darstellen liessen; es wurde
also nothwendig, die Differenzen so gut als möglich auf alle drei Beobachtungen zu vertheilen. Dies
habe ich auf's sorgfältigste gethan, und zwar nach einer eignen Methode, die ich auch in mein Werk
mit aufzunehmen denke. Auf diese Weise sind also folgende neue parabolische Elemente des zweiten
Cometen entstanden: Seeberger Zeit. Bewegung rechtläufig.\par}



\begin{center}\begingroup
\fontsize{9}{11}\selectfont
\setlength{\tabcolsep}{2pt}\renewcommand{\arraystretch}{1.12}
\begin{tabular}{lr}
Durchg.\ d.\ das Perihelium 1805 Dec.\ 31 & $6^{\mathrm{u}}\,54^{\mathrm{m}}\,34^{\mathrm{s}}$ \\
Länge des Periheliums .\ .\ .\ .\ .\ . & $109^{\circ}\,21'\,50''$ \\
Länge des aufsteigenden Knotens .\ . & $250\;33\;35$ \\
Neigung der Bahn .\ .\ .\ .\ .\ .\ . & $16^{\circ}\,30'\,32''$ \\
Logarithmus des kleinsten Abstandes .\ . & 9.9503300 \\
Logarithmus der mittl.\ tägl.\ Bewegung . & 0.0346333 \\
\end{tabular}\endgroup\end{center}



\par\indent Hier folgt nunmehr die Vergleichung dieser Elemente mit den sämmtlichen Beobachtungen:



\begin{center}\begingroup
\fontsize{8}{10}\selectfont
\setlength{\tabcolsep}{2pt}\renewcommand{\arraystretch}{1.12}
\begin{tabular}{|c|c|c|c|c|c|l|}
\hline
& & \shortstack{Berechnete\\gerade\\Aufsteigung} & \shortstack{Berechnete\\Abweichung} & \multicolumn{2}{c|}{Unterschied in der} & \\
1805 & Mittlere Zeit & & & \shortstack{ger.\ Aufsteig.\\R---B} & \shortstack{Abweichung\\R---B} & Beobachter \\
\hline
Nov.\ 15 & $10\,28^{\mathrm{m}}\,13^{\mathrm{s}}$ & $14^{\circ}\,33'\,0''$ & $39^{\circ}\,18'\,41''$ N. & $-2'\,30''$ & $+0'\,26''$ & Mr.\ Thulis \\
16 & $7\;\;17\;\;54$ & $14\;\;12\;\;7$ & $38\;\;58\;\;53$ & $-1\;\;58$ & $+0\;\;46$ & --- --- \\
19 & $7\;\;11\;\;17$ & $12\;\;55\;\;12$ & $37\;\;34\;\;27$ & $-0\;\;4$ & $+1\;\;7$ & --- --- \\
20 & $7\;\;3\;\;15$ & $12\;\;28\;\;23$ & $36\;\;59\;\;41$ & $-0\;\;23$ & $-0\;\;32$ & --- --- \\
21 & $9\;\;36\;\;1$ & $11\;\;57\;\;5$ & $36\;\;16\;\;16$ & $+1\;\;13$ & $-0\;\;20$ & --- --- \\
23 & $7\;\;23\;\;15$ & $11\;\;1\;\;15$ & $34\;\;46\;\;43$ & $+1\;\;0$ & $+0\;\;40$ & --- --- \\
29 & $6\;\;39\;\;44$ & $7\;\;24\;\;39$ & $26\;\;47\;\;49$ & $+1\;\;32$ & $+0\;\;1$ & --- --- \\
Dec.\ 2 & $5\;\;34\;\;9$ & $4\;\;48\;\;43$ & $18\;\;58\;\;18$ & $+1\;\;35$ & $-0\;\;52$ & Dr.\ Olbers \\
2 & $9\;\;39\;\;0$ & $4\;\;36\;\;3$ & $18\;\;20\;\;18$ & $+1\;\;47$ & $-0\;\;3$ & Mr.\ Thulis \\
3 & $5\;\;48\;\;39$ & $3\;\;39\;\;51$ & $15\;\;4\;\;52$ & $+0\;\;45$ & $-2\;\;5$ & Dr.\ Olbers \\
3 & $6\;\;30\;\;37$ & $3\;\;36\;\;52$ & $14\;\;55\;\;11$ & $+1\;\;10$ & $-0\;\;26$ & Mr.\ Thulis \\
4 & $6\;\;31\;\;5$ & $2\;\;16\;\;36$ & $10\;\;6\;\;29$ N. & $+2\;\;12$ & $-1\;\;12$ & --- --- \\
6 & $7\;\;16\;\;16$ & $358\;\;39\;\;22$ & $3\;\;39\;\;32$ S. & $+0\;\;22$ & $+0\;\;48$ & --- --- \\
7 & $7\;\;35\;\;21$ & $356\;\;9\;\;21$ & $13\;\;4\;\;56$ & $0\;\;0$ & $+0\;\;42$ & --- --- \\
8 & $5\;\;27\;\;22$ & $353\;\;20\;\;30$ & $22\;\;56\;\;24$ & $-1\;\;10$ & $+2\;\;44$ & Dr.\ Olbers \\
8 & $6\;\;46\;\;23$ & $353\;\;8\;\;33$ & $23\;\;34\;\;40$ & $-0\;\;42$ & $+2\;\;18$ & --- --- \\
8 & $6\;\;55\;\;17$ & $353\;\;7\;\;44$ & $23\;\;35\;\;37$ & $+0\;\;4$ & $-0\;\;47$ & Dr.\ Gauss \\
8 & & $353\;\;6\;\;21$ & $23\;\;40\;\;59$ & $-0\;\;33$ & $-0\;\;15$ & Dr.\ Maskelyne \\
8 & $6\;\;59\;\;39$ & $353\;\;3\;\;54$ & $23\;\;47\;\;34$ & $+0\;\;9$ & $-0\;\;33$ & Mr.\ Thulis \\
8 & $7\;\;58\;\;3$ & $352\;\;57\;\;34$ & $24\;\;6\;\;39$ & $-0\;\;23$ & .\quad . & Dr.\ Gauss \\
9 & $6\;\;4\;\;9$ & $349\;\;2\;\;44$ & $35\;\;23\;\;39$ & $-0\;\;16$ & .\quad . & Mr.\ Thulis \\
\hline
\end{tabular}\endgroup\end{center}



\par\indent Diese Vergleichung ist in aller Schärfe geführt, auf Aberration, Nutation und Parallaxe gehörig
\SourceLeaf{277}%
Rücksicht genommen, und die Theile von Secunden, die grösserer Genauigkeit wegen mit beobachtet
sind, erst am Ende der Rechnung weggelassen. [Die Zeit ist die an den Beobachtungsorten.]

\par\indent Obgleich die Differenzen an sich nicht viel grösser sind, als gewöhnlich Beobachtungs-Fehler bei
Cometen wol sein können, so zeigt sich doch in ihrem Gange eine unverkennbare Regelmässigkeit, die
ein Beweis für die Güte der \textsc{Thulis}'schen Beobachtungen zu sein scheint. In der That
schimmerte nach meiner eigenen Erfahrung vom 8. December der Kern so bestimmt durch, dass sich
dieser Comet mit ziemlicher Schärfe hat müssen beobachten lassen, und es ist sehr zu bedauern, dass
nicht mehrere Meridian-Beobachtungen bekannt sind, als die Greenwicher vom 8. December (aus den
Zeitungen) und die Marseiller vom 9\textsuperscript{ten}. Aus dieser Regelmässigkeit liess sich
voraus sehen, dass man durch eine ohne die parabolische Hypothese bestimmte Bahn, sie möchte nun
eine Ellipse oder eine Hyperbel werden, wenigstens eine bedeutend grössere Übereinstimmung mit den
Beobachtungen würde erreichen können. Ich habe daher diese Arbeit um so lieber übernommen, da ich
gerade um diese Zeit mit einigen wichtigen Abänderungen und Zusätzen zu meiner Methode beschäftigt
war, wovon ich eine practische Anwendung zu haben wünschte, zumal da eben \emph{dieses} Beispiel
etwas Eigenthümliches hat, was bei den Asteroiden Ceres, Pallas und Juno, nicht vorkommen kann, und
worüber hier nicht der Ort ist, mich weitläuftiger zu erklären.

\par\indent Das Resultat dieser Arbeit ist wirklich von mehr als einer Seite merkwürdig genug. Ich erhielt
nämlich folgende \emph{elliptische} Bahn:

\begin{center}\begingroup
\fontsize{9}{11}\selectfont
\setlength{\tabcolsep}{2pt}\renewcommand{\arraystretch}{1.12}
\begin{tabular}{lr}
\multicolumn{2}{l}{Durchgang durch das Perihelium} \\
\quad den 2. Januar 1806 Seeb. Zeit .\ . & $11^{\mathrm{u}}\;8^{\mathrm{m}}\;45^{\mathrm{s}}$ \\
Länge des Periheliums .\ .\ .\ .\ . & $109^\circ\;30'\;2''3$ \\
Länge des aufsteigend. Knotens .\ . & $251\;28\;22.5$ \\
Neigung der Bahn .\ .\ .\ .\ .\ .\ . & $12\;43\;10.0$ \\
Logarithm. d. halben grossen Axe .\ .\ . & $0.4505887$ \\
Tägliche mittlere Bewegung .\ .\ .\ .\ . & $748''383$ \\
Umlaufszeit .\ .\ .\ .\ .\ .\ .\ . & 1731 Tage 17 St. \\
Excentricität .\ .\ .\ .\ .\ .\ .\ .\ .\ . & $0.6769242$ \\
Grösste Distanz von der Sonne .\ .\ .\ . & $4.732625$ \\
Kleinste \quad---\ ---\ --- .\ .\ .\ . & $0.911786$ \\
Logarithmus der kleinsten Distanz .\ .\ . & $9.9598931$ \\
\end{tabular}\endgroup\end{center}

\Needspace{28\baselineskip}
\par\indent Diese Elemente stimmen mit den Beobachtungen folgendermaassen:

\begin{center}\begingroup
\fontsize{9}{11}\selectfont
\setlength{\tabcolsep}{2pt}\renewcommand{\arraystretch}{1.12}
\begin{tabular}{c|cc|cc|l}
& \shortstack{Berechnete gerade\\Aufsteigung} & \shortstack{Berechnete\\Abweichung} & \multicolumn{2}{c|}{Unterschied in der} & \\
1805 & & & \shortstack{ger.\ Aufsteig.\\R---B} & \shortstack{Abweichung\\R---B} & Beobachter \\
\hline
Nov. 15 & $14^\circ\;35'\;45''$ & $39^\circ\;17'\;56''$ N. & $+15''$ & $-19''$ & Mr. Thulis \\
16 & $14\;14\;20$ & $38\;56\;6$ & $+15$ & $-1$ & --- --- \\
19 & $12\;55\;23$ & $37\;34\;5$ & $+7$ & $+45$ & --- --- \\
20 & $12\;27\;59$ & $36\;59\;24$ & $-47$ & $-49$ & --- --- \\
21 & $11\;55\;59$ & $36\;15\;58$ & $+7$ & $-38$ & --- --- \\
23 & $10\;59\;33$ & $34\;46\;25$ & $-42$ & $+22$ & --- --- \\
29 & $7\;22\;34$ & $26\;47\;40$ & $-33$ & $-8$ & --- --- \\
Dec.\ 2 & $4\;47\;28$ & $18\;58\;28$ & $+20$ & $-42$ & Dr. Olbers \\
2 & $4\;34\;22$ & $18\;20\;44$ & $+6$ & $+13$ & Mr. Thulis \\
3 & $3\;38\;51$ & $15\;5\;11$ & $-15$ & $-1'\;46$ & Dr. Olbers \\
3 & $1\;35\;50$ & $14\;55\;37$ & $+8$ & $0$ & Mr. Thulis \\
4 & $2\;15\;55$ & $10\;7\;4$ N. & $+1'\;31$ & $-37$ & --- --- \\
6 & $358\;39\;13$ & $3\;38\;58$ S. & $+13$ & $+14$ & --- --- \\
7 & $356\;9\;28$ & $13\;4\;43$ & $+7$ & $+29$ & --- --- \\
8 & $353\;21\;10$ & $22\;56\;53$ & $-30$ & $+3\;13$ & Dr. Olbers \\
8 & $353\;9\;12$ & $23\;35\;19$ & $-3$ & $+2\;57$ & --- --- \\
8 & $353\;8\;10$ & $23\;36\;16$ & $+30$ & $-8$ & Dr. Gauss \\
8 & $353\;6\;55$ & $23\;41\;36$ & $+1$ & $+22$ & Dr. Maskelyne \\
8 & $353\;4\;18$ & $23\;48\;8$ & $+33$ & $+1$ & Mr. Thulis \\
8 & $352\;57\;46$ & $24\;7\;18$ & $-11$ & .\quad . & Dr. Gauss \\
9 & $349\;2\;56$ & $35\;25\;21$ & $-4$ & .\quad . & Mr. Thulis \\
\end{tabular}\endgroup\end{center}
\SourceLeaf{278}%
\par\indent Bei der Vergleichung dieser Differenzen mit den obigen in der Parabel gefundenen
zeigt sich auf den ersten Blick, dass die Ellipse um \emph{sehr vieles} besser
harmonirt. Unter \textsc{Thulis} Beobachtungen findet sich lediglich nur die
gerade Aufsteigung vom 4. December, wo die Differenz über $50''$ geht, und an
diesem Tage war ein Stern gebraucht, den ich aus der Conn. des tems, An. XIII,
genommen habe, von dem ich aber in der Hist. cél. keine Beobachtungen finden
konnte. Vielleicht bedarf dieser einer Berichtigung. Dr. \textsc{Olbers} hatte
selbst seine Declination vom 3. December als sehr zweifelhaft angegeben, und am
8. December konnte man die Rechnung mit Dr. \textsc{Olbers} Declinationen nicht
in Übereinstimmung bringen, ohne sich von der Greenwicher Meridian-Beobachtung,
womit auch \textsc{Thulis} und ich bis auf Kleinigkeiten harmoniren, um 3
Minuten zu entfernen. Auch die Declination am 9. December stimmt, wenn man sie
durch Refraction auf scheinbare bringt, mit \textsc{Thulis} Schätzung besser in
der Ellipse, als in der Parabel.

\par\indent Hieraus lassen sich nun folgende Folgerungen ziehen:

\par\indent Die Beobachtungen des Cometen lassen sich durch alle Ellipsen, die von der
obigen den Übergang zur Parabel machen, oder deutlicher, sie lassen sich in der
Ellipse mit \emph{jeder} mittlern Entfernung, die nur grösser als 2.82 ist,
immer noch besser darstellen, als in irgend einer Parabel, wiewol successive
immer weniger gut, so wie sie sich mehr von obiger Ellipse entfernt. Auch auf
der andern Seite wird man mit einer noch kleinern mittlern Entfernung als 2.82
bis auf eine gewisse Gränze die Beobachtungen noch besser darstellen können,
als in der Parabel, da die Differenzen zu beiden Seiten des Minimum nur nach
und nach anwachsen. Ich habe zwar hierüber keine Rechnung angestellt, glaube
aber doch, dass selbst die mittlere Distanz 2.0 oder eine Umlaufszeit von
weniger als drei Jahren sich noch eben so gut mit allen Beobachtungen
vertrage, als die Parabel.

\par\indent Ich habe schon öfters gedacht und geäussert, dass unter den hundert bisher
berechneten Cometen, wo man immer ohne weiteres die parabolische Hypothese
anwendet, und schon zufrieden ist, wenn sich diese mit den Beobachtungen
verträgt, bei weitem der grössere Theil wol von der Art sein möchte, dass sich
aus den Beobachtungen allein noch nicht einmal gewiss beweisen liesse, dass
ihre Bahn der Parabel sehr nahe kommen müsste. Es wäre daher zu wünschen, dass
man bei jedem berechneten Cometen eine Idee davon hätte, wie weit die
Schranken von einander liegen, innerhalb deren die Bahn liegen muss, wenn die
Rechnung sich nicht zu sehr von der Beobachtung entfernen soll. In dieser
Rücksicht ist meiner Meinung nach unser Comet, der in 24 Tagen eine
geocentrische Bewegung von beinahe 80 Graden gezeigt hat, ein sehr
merkwürdiges Beispiel, wie sehr wir über die wahren Dimensionen der Bahn in
Ungewissheit bleiben. Indessen gestehe ich gern, dasss wenn der Comet noch bis in den Januar hinein gut beobachtet worden wäre (was
in der südlichen Hemisphäre sehr gut hätte geschehen können), sich schon
etwas ziemlich Entschiedenes über seine wahre Bahn müsste sagen lassen. Es
würde mir in dieser Hinsicht schon willkommen sein, wenn Sie mir
\textsc{Thulis} Original-Beobachtungen vom 10. und 12. November verschaffen
könnten, so wie, früh oder spät, eine genaue Bestimmung der am 11. November
verglichenen Sterne. Eine um fünf Tage verlängerte Dauer der Beobachtung hat
hier schon einigen Werth.

\par\indent Bei aller dieser Ungewissheit schien doch die um so vieles bessere
Übereinstimmung der obigen Ellipse es einigermaassen wahrscheinlich zu
machen, dass sich dieser Comet wirklich in einer Ellipse von einer nicht gar
zu grossen Umlaufszeit bewege. Hiezu kam noch die Ähnlichkeit mit dem Cometen
vom Jahre 1772. \textsc{Bessel} hatte zwar durch seine Rechnungen gefunden,
dass, wenn man jeden dieser Cometen in einer Ellipse von 33 Jahren
Umlaufszeit darstelle, die übrigen Elemente doch eben keine grössere
Übereinstimmung zeigen, als bei den Parabeln Statt gefunden hatte. Indessen
blieb es nach obiger Untersuchung immer noch möglich, dass ein und derselbe
Comet in diesen 33 Jahren vielleicht zwei, drei
\SourceLeaf{279}%
oder mehrere Umläufe gemacht hätte. Die Lichtschwäche dieses im vorigen Jahre nur durch seine grosse
Nähe so augenfällig gewordenen Cometen würde es leicht erklären, dass er bei den übrigen Durchgängen
durch's Perihelium, die vielleicht unter ungünstigen Umständen geschahen, nicht bemerkt wurde.
Allein meine hierüber angestellten Rechnungen ergeben, dass auch dieses nicht angehen wird. Ich finde
nemlich, dass man den Beobachtungen der zwei Cometen zugleich einerlei kleinsten und grössten Abstand
von der Sonne nicht anpassen kann, ohne bei den übrigen Elementen (als Länge der Sonnennähe,
Neigung der Bahn und Länge des Knotens) eine zu bedeutende Differenz zu erhalten.

\par\indent Hier ist das Tableau von zweierlei Elementen des Cometen von 1772; erstens in einer Parabel,
wo ich den kleinsten Abstand so gross genommen hatte, als er in meinen (frühern) parabolischen Elementen
des vorjährigen Cometen war; zweitens in einer Ellipse von derselben kleinsten und grössten
Distanz, wie in der obigen. Die respective Vergleichung dieser Elemente mit denen für den vorjährigen
Cometen zeigt, dass man bei dem Cometen von 1772, unter der Voraussetzung, dass die grösste und
kleinste Entfernung dieselbe wie bei dem von 1805 sei, immer die Länge der Sonnennähe ansehnlich
kleiner, die Länge des Knotens ansehnlich grösser, und auch die Neigung $4^\circ$ bis $5^\circ$ grösser findet, als
bei letzterem.

\begin{center}\begingroup
\fontsize{8}{10}\selectfont
\setlength{\tabcolsep}{2pt}\renewcommand{\arraystretch}{1.12}
\begin{tabular}{l|c|c}
& \multicolumn{2}{c}{Elemente des Cometen vom Jahre 1772.} \\
 & parabolische & elliptische \\
\hline
Durchgang durch das Perihelium & d. 9. Febr. 1772. $5^{\mathrm{u}}$ & d. 8. Febr. 1772. $1^{\mathrm{u}}$ \\
Länge des Periheliums . . . . & $90^{\circ}\,17'$ & $97^{\circ}\,21'$ \\
Länge des aufsteigenden Knotens & $261\ \ 9$ & $263\ \ 24$ \\
Neigung der Bahn . . . . . & $20\ \ 28$ & $17\ \ 39$ \\
kleinster Abstand . . . . . . & $0.8918$ & $0.9118$ \\
grösster Abstand . . . . . . & $\infty$ & $4.7326$ \\
\end{tabular}\endgroup\end{center}

\par\indent Diese Elemente gründen sich nur auf die beiden äussern \textsc{Messier}'schen Beobachtungen (zwei Beobachtungen
sind immer hinreichend, wenn zwei Stücke der Elemente als gegeben angesehen werden);
und dann ergab sich bei der mittlern Beobachtung eine Differenz von 6 Minuten bei der Parabel, und
von 3 Min. bei der Ellipse in der Länge; die Breite stimmte bei beiden gut. Durch Vertheilung würden
sich indess diese Differenzen etwa auf die Hälfte reduciren lassen, ohne dass dadurch die Elemente selbst
\emph{stark} geändert werden könnten.

\par\indent Dessen ungeachtet will ich jetzt noch nicht über die Möglichkeit der Identität der beiden Cometen
absprechen. Die bisherigen Untersuchungen zeigen allerdings, dass sich beider Bewegungen nicht
befriedigend durch einerlei Bahn darstellen lassen; allein noch bleibt die Möglichkeit übrig, dass sich
die Elemente durch eine fremdartige Störung ansehnlich geändert haben. Zwar findet sich, wenn man
dem Cometen von 1772 eine Umlaufszeit von 33 Jahren gibt, so dass er 1805 wieder in's Perihelium
kommen musste, nicht, dass er einem Hauptplaneten nahe gekommen wäre; allein noch bleibt die Möglichkeit
übrig, dass sowohl der Ellipse, aus deren Perihelium er 1772 ging, als der, in deren Perihelium er
1805 ankam, ganz verschiedene Umlaufszeiten zugehören, und dass er irgend einmal in der Zwischenzeit
einem Hauptplaneten so nahe gekommen ist, dass dieser die erste Ellipse in die zweite verwandeln konnte,
es sei nun, dass er in dieser Zeit einen oder mehrere Umläufe gemacht hat. Es bleibt also noch die
freilich etwas weitläuftige Untersuchung übrig, ob unter allen den elliptischen Bahnen, wodurch sich die
Bewegungen der beiden Cometen darstellen lassen (und die, wie wir bestimmt wissen, in so weite Gränzen
eingeschlossen sind), sich nicht zwei von der Art finden, dass jeder Comet durch die seinige für irgend
einen Tag zwischen 1772 und 1805 einen und denselben Ort erhalten, bei welchem sich an eben dem Tage
ein Hauptplanet sehr nahe befand. In einem solchen Falle musste dieser Hauptplanet einen starken Ein%
\SourceLeaf{280}%
fluss auf die Elemente haben, der aber sehr verschieden ausfallen musste, je nachdem die relative Bewegung so oder anders war; woraus eine, wie auch immer beschaffene Änderung der Elemente leicht erklärlich gemacht werden könnte. Hierüber denke ich künftig noch Untersuchungen anzustellen. --- --- ---
